\origpage{240}{ROBERTSON AND SEYMOUR}
We claim that $(T, \tau')$ is rotund. For let $P$ be a directed path of $T$ with first edge $f_2$ and last edge $f_1$, such that $f_1, f_2$ have the same order $k$ and every other edge of $P$ has order $\ge k$ in $(T, \tau')$. Let the first and last vertices of $P$ be $t_2, t_1$. By (2), there are $k$ mutually vertex-disjoint paths $Q_1, \dots, Q_k$ of $K$ between $V(\tau(t_1))$ and $V(\tau(t_2))$. Let $(H_1, H_2)$ be a separation of $G^-$ with $(\tau' \times T^{f_1})^- \subseteq H_1$ and $(\tau' \times T_{f_2})^- \subseteq H_2$. Each $Q_i$ has one end in $V(H_1)$ and the other end in $V(H_2)$, and for each $e \in E(Q_i)$ either both ends of $e$ in $K$ are ends in $G$ of some $e' \in E(G)$ or both ends of $e$ in $K$ are in $\bar{\pi}(G)$; and in either case both ends of $e$ in $K$ lie in some $V(\tau'(t))$ and hence either both lie in $V(H_1)$ or both lie in $V(H_2)$. Hence each $V(Q_i)$ meets $V(H_1 \cap H_2)$, and so $|V(H_1 \cap H_2)| \ge k$. Thus $(H_1, H_2)$ has order $\ge k$; and hence $(T, \tau')$ is rotund, as required. $\square$

We introduce a fourth axiom. For each $n \ge 0$, let $\mathscr{S}_n$ denote the class of all star-decompositions $(S, \sigma)$ with $\sigma \times S \in \mathscr{R}$ and with $|V(\sigma(o(S)))| \le n$.

\textit{Axiom} 4. \textit{For each $n \ge 0$, $\mathscr{S}_n$ is good.}

From (5.2) we deduce

\stmt
(5.3) \textit{Let $G_i$ ($i = 1, 2, \dots$) be a countable sequence of elements of $\mathscr{R}$, each of tree-width $\le n$. Then there exist $i' > i \ge 1$ such that $G_i$ is simulated in $G_{i'}$.}

\textit{Proof.} From (5.2), for each $i \ge 1$ there is a rotund tree-decomposition of $G_i$ of width $\le n$, and hence with all its branchings in $\mathscr{S}_{n + 1}$. Since $\mathscr{S}_{n + 1}$ is good and symmetric, the result follows from (4.2). $\square$

\textit{Proof of} (1.5) (\textit{Sketch}). Let $G_i$ ($i = 1, 2, \dots$) be a countable sequence of graphs, each of tree-width $\le n$. Let $\mathscr{R}$ be the class of all rooted graphs $G$ such that $G^-$ is a subgraph of some $G_i$. For $H_1, H_2 \in \mathscr{R}$ we say that $H_1$ is simulated in $H_2$ if $|\bar{\pi}(H_1)| = |\bar{\pi}(H_2)|$ and there exists $H \in \mathscr{R}$ with $H^- \subseteq H_2^-$ and $\pi(H) = \pi(H_2)$, such that a rooted graph isomorphic to $H_1$ can be obtained from $H$ by edge-contraction. We verify Axioms 1--4, and (1.5) follows from (5.3). $\square$

We have omitted verifying Axioms 1--4, because it is easy and we shall later carry out the verification for a more general ``concrete'' definition of simulation. \fix{This verification is carried out in Sections 6--9, and a complete proof of (1.5), by way of
(9.1), is given on page 253.\fixnote{fix:1.5-sketch}{The paper does not verify Axioms 1--4 for the
relation defined in the sketch above. The derivation of (1.5) from (9.1) on page 253 is complete,
and the details added there show that simulation of the patchworks used gives minors.}}

\sectionhead{6. \textsc{Patchworks}}

Now we come to the second part of the paper. We introduce a concrete containment relation on a class of rooted hypergraphs, and verify that it satisfies the axioms.
